\chapter{Design and Implementation Specifications}

This chapter describes the architecture, fault injection mechanism, simulator design, and implementation details of the artefact.

\section{System Architecture}

The artefact is structured as an AWS Serverless Application Model (SAM) application with two arms: a synchronous API control and an SQS-based queue treatment. Both arms process synthetic order payloads, implement idempotency, emit metrics, and support fault injection.

\subsection{AWS SAM Stack (\texttt{template.yaml})}

The SAM template defines the following resources:

\begin{description}
    \item[SQS Primary Queue (\texttt{OrdersQueue}):] Standard queue with configurable \texttt{VisibilityTimeout}, \texttt{DelaySeconds}, and \texttt{RedrivePolicy} pointing to the DLQ with \texttt{maxReceiveCount}.
    
    \item[SQS Dead-Letter Queue (\texttt{OrdersDLQ}):] Standard queue for failed messages. Messages moved here after exceeding \texttt{maxReceiveCount} retries.
    
    \item[Lambda Queue Consumer (\texttt{QueueConsumerFunction}):] Python 3.12 function triggered by SQS event source mapping with configurable \texttt{BatchSize}, \texttt{MaximumBatchingWindowInSeconds}, and \texttt{FunctionResponseTypes: [ReportBatchItemFailures]} for partial batch handling.
    
    \item[Lambda Sync API (\texttt{SyncApiFunction}):] HTTP API Gateway + Lambda for synchronous control arm.
    
    \item[DynamoDB Table (\texttt{OrdersTable}):] Stores processed orders with \texttt{order\_id} as partition key. Used for idempotency checks and result persistence.
    
    \item[SSM Parameter (\texttt{FaultSwitchParameter}):] Stores fault injection configuration as JSON. Both Lambda functions read this to determine fault behavior.
    
    \item[CloudWatch Metrics:] Custom namespace \texttt{/sqs-rr/metrics} emits \texttt{OrdersSent}, \texttt{OrdersProcessed}, \texttt{OrdersFailed}, \texttt{DLQDepth}, \texttt{Latency}.
\end{description}

\subsection{Two-Arm Design}

\textbf{Sync Arm:} Producer sends POST requests to API Gateway → Lambda processes synchronously, checks idempotency in DynamoDB, applies fault if active, stores result, returns HTTP response.

\textbf{Queue Arm:} Producer sends messages to SQS → Lambda triggered by event source mapping with batch → processes each message, checks idempotency, applies fault, stores result or returns \texttt{batchItemFailures} → SQS redelivers failed messages or moves to DLQ after \texttt{maxReceiveCount}.

\section{Fault Injection Mechanism}

The fault injection switch is implemented as an in-process module (\texttt{src/common/injected\_fault\_handler.py}) integrated into both Lambda handlers.

\subsection{Fault Types}

\begin{description}
    \item[\texttt{consumer\_kill}:] Consumer is unavailable (simulated by returning immediately without processing). In live AWS, this would be a Lambda container crash or timeout; in the simulator, messages sit in-flight until visibility timeout expires.
    
    \item[\texttt{unhandled\_error}:] Raises \texttt{RuntimeError} during processing. In Lambda, this triggers a retry (redelivery); in simulator, same effect.
    
    \item[\texttt{datastore\_reject}:] DynamoDB \texttt{put\_item} fails with \texttt{ConditionalCheckFailedException} (simulated rejection). Tests partial failure reporting.
    
    \item[\texttt{datastore\_timeout}:] DynamoDB \texttt{put\_item} is delayed (3-5 seconds in simulator), simulating overload. Tests interaction with visibility timeout.
\end{description}

\subsection{Fault Controller}

The controller reads fault config from SSM parameter (or simulator config) structured as:
\begin{verbatim}
{
  "fault_type": "consumer_kill",
  "start_time": 30,
  "stop_time": 40,
  "probability": 1.0
}
\end{verbatim}

At each invocation, the handler checks current time against \texttt{start\_time} and \texttt{stop\_time}. If within fault window and probability threshold met, the fault is injected.

\section{Local Simulator}

The simulator (\texttt{src/localsim/engine.py}) implements discrete-event semantics for SQS, Lambda, and DynamoDB without AWS network calls.

\subsection{Discrete-Event Engine}

\begin{itemize}
    \item \textbf{Virtual Clock:} All events are scheduled on a priority queue ordered by timestamp. The engine advances time by processing the next event.
    
    \item \textbf{Event Types:} \texttt{MessageSent}, \texttt{MessageVisible}, \texttt{LambdaPoll}, \texttt{LambdaInvoke}, \texttt{MessageDeleted}, \texttt{MessageDLQ}, \texttt{MetricsSnapshot}.
    
    \item \textbf{Message States:} \texttt{pending} (sent but delayed), \texttt{visible} (available for receive), \texttt{inflight} (received, visibility timeout active), \texttt{deleted} (processed), \texttt{dlq} (dead-lettered).
\end{itemize}

\subsection{SQS Simulation (\texttt{src/localsim/sqs.py})}

\begin{itemize}
    \item \texttt{send\_message(order, delay)}: Enqueues message, schedules \texttt{MessageVisible} event after delay.
    
    \item \texttt{receive\_messages(batch\_size)}: Returns up to \texttt{batch\_size} visible messages, transitions to \texttt{inflight}, schedules visibility timeout expiry.
    
    \item \texttt{delete\_message(receipt)}: Removes message from queue, cancels visibility timeout event.
    
    \item \texttt{visibility\_timeout\_expired(receipt)}: Returns message to \texttt{visible} state, increments \texttt{receive\_count}. If \texttt{receive\_count >= maxReceiveCount}, moves to DLQ.
\end{itemize}

\subsection{Lambda Event Source Mapping}

The simulator polls the SQS queue every \texttt{poll\_interval} (default 1 second). When messages are available, it batches them and invokes the handler (\texttt{src/queue\_consumer/handler.py}) via Python function call (not AWS invocation). The handler returns \texttt{batchItemFailures}; simulator deletes successes and leaves failures in-flight.

\subsection{Idempotency Layer (\texttt{src/common/idempotency.py})}

Stores processed \texttt{order\_id} in DynamoDB (or in-memory dict for simulator). Before processing, checks if \texttt{order\_id} exists; if yes, skips processing (idempotent). This prevents duplicate stores for redelivered messages.

\section{Experiment Runner (\texttt{src/control/experiment\_runner.py})}

The runner walks a config matrix, generates run manifests, executes each run, aggregates metrics, and persists results.

\subsection{Config Matrix}

YAML files (\texttt{configs/*.yaml}) define experiment phases:
\begin{verbatim}
base:
  visibility_timeout: [30, 60, 120, 300, 600, 900]
  max_receive_count: [1, 3, 5, 10]
  batch_size: [1, 10, 25, 50]
  delivery_delay: [0, 5, 60, 300]
campaigns:
  - name: consumer_kill
    fault_type: consumer_kill
    fault_window: [30, 40]
repeats: 5
\end{verbatim}

The runner generates all combinations, randomizes order (seeded), and assigns unique run IDs (hash of config + seed).

\subsection{Manifest Structure}

Each run produces a JSON manifest:
\begin{verbatim}
{
  "run_id": "baseline_vt300_batch10_rep3",
  "config": {...},
  "seed": 42,
  "metrics": {
    "loss": 0.0,
    "duplicates": 0.002,
    "dlq_rate": 0.05,
    "recovery_time_backlog": 35.2,
    "throughput": 287.4,
    "latency_p50": 0.12
  },
  "backend": "localsim",
  "timestamp": "2026-09-19T08:30:00Z"
}
\end{verbatim}

Manifests are committed to \texttt{results/<phase>/manifests/<run\_id>.json}.

\section{Metrics Pipeline (\texttt{src/control/collect\_metrics.py})}

\subsection{Loss Calculation}

\[
\text{loss} = 1 - \frac{\text{unique\_stored} + \text{dlq\_count}}{\text{sent\_count}}
\]

\subsection{Recovery Time}

\begin{enumerate}
    \item Record \texttt{baseline\_depth} (queue depth at \texttt{t = fault\_start - 5}).
    \item Measure \texttt{depth(t)} for \texttt{t > fault\_stop}.
    \item Recovery time = first \texttt{t} where \texttt{depth(t) <= baseline\_depth}.
\end{enumerate}

Separate metrics for \texttt{total\_backlog} (visible + inflight) and \texttt{visible\_only}.

\subsection{Duplicate Rate}

\[
\text{duplicates} = \frac{\text{total\_stored} - \text{unique\_stored}}{\text{unique\_stored}}
\]

\section{Statistical Analysis Scripts}

\texttt{analysis/stats\_tests.py} implements:
\begin{itemize}
    \item Hypothesis tests (Mann-Whitney U, Kruskal-Wallis H)
    \item Effect size (rank-biserial r, epsilon-squared η²)
    \item Holm-Bonferroni correction
    \item Bootstrap 95\% CIs
    \item Power analysis
\end{itemize}

\texttt{analysis/plot\_results.py} generates all figures using matplotlib/seaborn.

\section{Testing and Quality Gates}

\subsection{Unit Tests (\texttt{tests/unit/})}

\begin{itemize}
    \item \texttt{test\_faults.py}: Fault injection logic
    \item \texttt{test\_sqs\_sim.py}: SQS queue semantics (visibility timeout, receive count, DLQ)
    \item \texttt{test\_engine.py}: Discrete-event engine ordering
    \item \texttt{test\_handler.py}: Lambda handler with mocked SQS
    \item \texttt{test\_metrics.py}: Metrics aggregation
    \item \texttt{test\_idempotency.py}: Deduplication logic
    \item \texttt{test\_cost.py}: Cost estimation
    \item \texttt{test\_config.py}: Config matrix generation
\end{itemize}

\subsection{Integration Tests (\texttt{tests/integration/})}

\begin{itemize}
    \item \texttt{test\_aws\_moto.py}: SAM stack deployment with moto (mocked AWS)
    \item \texttt{test\_dry\_run\_pipeline.py}: End-to-end pilot run in simulator
\end{itemize}

\subsection{CI/CD}

GitHub Actions (\texttt{.github/workflows/main.yml}) runs \texttt{pytest}, \texttt{cfn-lint}, and \texttt{make gates} on every push.

\section{Repository Structure}

\begin{verbatim}
sqs-reliability-recovery/
├── template.yaml           SAM stack
├── src/
│   ├── queue_consumer/     Lambda handler (queue arm)
│   ├── sync_api/           Lambda handler (sync arm)
│   ├── common/             Fault injection, idempotency, metrics
│   ├── producer/           Synthetic order generator
│   ├── control/            Experiment runner, manifests, cost guard
│   └── localsim/           Simulator
├── configs/                Experiment phase definitions
├── analysis/               stats_tests.py, plot_results.py
├── results/                Manifests, summaries, figures
├── tests/                  Unit + integration
├── docs/                   Config manual, architecture, demo
├── Makefile                Build + test + run shortcuts
├── requirements.txt        Dependencies
└── STATUS.md               Simulation vs AWS disclosure
\end{verbatim}

\section{Reproducibility}

All runs use seeded RNG (\texttt{random.seed(run\_seed)}). Run IDs are deterministic (hash of config). Manifests include seeds, timestamps, git SHAs. Commands:
\begin{verbatim}
make setup      # Create venv, install deps
make test       # Unit + integration tests
make pilot      # 18 pilot runs
make experiments # All phases (~350 runs, 30s)
make stats      # Hypothesis tests
make figures    # 11 plots
\end{verbatim}

\section{Limitations}

\begin{itemize}
    \item No cold start modeling (assumes warm Lambda).
    \item Deterministic timing (no network jitter).
    \item No concurrency limits (real AWS soft limit 1,000).
    \item No AWS service failures (assumes SQS/Lambda/DynamoDB never fail).
\end{itemize}

See \texttt{STATUS.md} for full disclosure.
